\begin{enumerate}[start=9]
	\item L'activité massique de l'échantillon radioactif peut être
	modélisée par la fonction $A$ définie sur l'intervalle $[0 ; +\infty[$ par :
	\[A(t) =5 \e^{-0,315t}\]
	Le temps $t$ est exprimé en année et l'activité massique $A(t)$ est exprimée en
	milliers de becquerel par kilogramme.

	On admet que $A$ est dérivable sur l'intervalle $[0 ; +\infty[ $ et on note $A'$ sa fonction
	dérivée.
	\begin{enumerate}
		\item  Vérifier que $A'(t) = -1,575 \e ^{-0,315t}$,
		\item  En déduire le sens de variation de la fonction $A$ sur l'intervalle $[0 ; +\infty[$.
	\end{enumerate}
	\item Résoudre l'équation $A(t) = 1,25$ dans l'intervalle $[0 ; +\infty[$.

	On donnera une valeur arrondie à $10^{-1}$.

	\item En déduire le temps nécessaire pour que l'échantillon de viande de Bécasse atteigne la limite de \np[Bq\cdot kg^{-1}]{1250} fixée par EURATOM.
\end{enumerate}
